\chapter{Botulism}

\section*{Definition and Overview}
Botulism is a rare but serious paralytic illness caused by botulinum toxin, one of the most potent poisons known. The toxin is produced by the bacterium \textit{Clostridium botulinum}, and in rare cases by related species such as \textit{C. baratii} and \textit{C. butyricum}. Botulinum toxin blocks the release of acetylcholine at the junction between nerves and muscles, preventing muscles from contracting and producing a characteristic descending paralysis that begins in the face and eyes and can progress to life-threatening respiratory failure. Botulism is not contagious between people. There are several forms: foodborne botulism follows ingestion of contaminated food; infant botulism occurs when spores germinate in the immature gut of babies; wound botulism follows contamination of a wound; and very rare iatrogenic or inhalational forms follow botulinum toxin injection or inhalation. All forms are medical emergencies requiring urgent hospital treatment.

\section*{Epidemiology}
Foodborne botulism is rare in Australia and other developed countries, with only a handful of cases per year, almost always linked to improperly processed home-preserved foods. Infant botulism is now the most common form in many countries and typically affects babies between two weeks and six months of age. Wound botulism is uncommon but has been increasingly recognised in people who inject drugs, particularly heroin. The illness occurs worldwide, with higher rates in regions where home canning and traditional fermented or smoked foods are common. All cases are notifiable to public health authorities so that contaminated food can be traced and recalled quickly.

\section*{Aetiology and Pathophysiology}
\begin{itemize}
    \item \textbf{Foodborne botulism:} spores survive in improperly processed home-preserved foods -- low-acid canned vegetables, fermented foods, cured meats, and vacuum-packed fish -- and germinate to produce toxin in the low-oxygen conditions inside the container.
    \item \textbf{Infant botulism:} spores from soil, dust, or occasionally honey germinate in the immature gut of babies under one year of age and release toxin in the intestine.
    \item \textbf{Wound botulism:} spores enter a deep or contaminated wound and produce toxin in the damaged, poorly oxygenated tissue.
    \item \textbf{Iatrogenic and inhalational forms:} rare, and follow injection of botulinum toxin for medical or cosmetic use or inhalation of aerosolised toxin.
    \item \textbf{Mechanism of action:} the toxin binds to nerve endings at the neuromuscular junction and cleaves proteins required for the release of acetylcholine. Without this chemical messenger, muscles never receive the signal to contract, producing flaccid paralysis.
    \item \textbf{Adult intestinal colonisation:} in rare cases, older children and adults with abnormal gut flora develop intestinal botulism similar to the infant form.
\end{itemize}

\section*{Clinical Features}
Symptoms usually begin 12 to 72 hours after eating contaminated food, although onset can be later, and the classic picture is a descending paralysis:
\begin{itemize}
    \item Double vision, blurred vision, and drooping eyelids (ptosis)
    \item Difficulty swallowing, speaking, and chewing
    \item A dry mouth, poor gag reflex, and sometimes a thick-feeling tongue
    \item Facial weakness with a blank, expressionless face
    \item Weakness spreading symmetrically down the arms, trunk, and legs
    \item Shortness of breath as the chest muscles weaken
    \item Constipation, often early and sometimes the presenting complaint
    \item In infants: poor feeding, a weak cry, a floppy "rag doll" posture, listlessness, and constipation
    \item Fever is usually absent, which helps distinguish botulism from many infections
\end{itemize}

\section*{Differential Diagnosis}
Other conditions cause similar acute weakness and must be distinguished quickly:
\begin{itemize}
    \item \textbf{Myasthenia gravis:} an autoimmune disease causing fatigable muscle weakness, often with similar eye and swallowing signs.
    \item \textbf{Guillain--Barr\'e syndrome:} including the Miller Fisher variant with eye movement problems, weakness, and unsteadiness.
    \item \textbf{Stroke or brainstem disorders:} sudden cranial nerve and motor deficits.
    \item \textbf{Poliomyelitis and other viral infections:} of the nervous system causing flaccid weakness.
    \item \textbf{Tick paralysis:} a rapidly progressive ascending weakness following a tick bite.
    \item \textbf{Chemical poisoning:} organophosphates and certain drugs can mimic the weakness and dry mouth.
    \item \textbf{Diphtheria:} can cause swallowing and breathing difficulty with a membrane in the throat.
\end{itemize}

\section*{When to Seek Medical Care}
Botulism is a medical emergency. Anyone with sudden double vision, drooping eyelids, difficulty swallowing or speaking, or progressive weakness should go to an emergency department immediately, especially if they have recently eaten home-preserved or suspect food. Do not wait to see whether symptoms resolve. A baby who is constipated, feeding poorly, floppy, or unusually quiet should also be assessed urgently.

\textbf{Call 000 (or local emergency services) for:}
\begin{itemize}
    \item Difficulty breathing or a feeling that breathing is becoming harder
    \item Difficulty swallowing, drooling, or choking
    \item Sudden weakness spreading from the face to the arms, trunk, or legs
    \item A baby who is floppy, feeding poorly, or breathing unusually
\end{itemize}

\section*{Diagnosis and Investigations}
Botulism is suspected on clinical grounds, because rapid treatment matters, and it is confirmed with laboratory tests:
\begin{itemize}
    \item \textbf{Clinical examination:} the distinctive pattern of descending, symmetrical, flaccid weakness with cranial nerve involvement and clear consciousness.
    \item \textbf{Toxin testing:} detection of botulinum toxin in serum, stool, or wound material, using a mouse bioassay or modern rapid methods.
    \item \textbf{Culture:} growth of \textit{Clostridium botulinum} from stool or food samples, which can take several days.
    \item \textbf{Testing of suspected food:} leftover food can be examined for toxin and the organism and linked to the patient by public health laboratories.
    \item \textbf{Nerve conduction studies and electromyography:} characteristic findings help distinguish botulism from Guillain--Barr\'e syndrome and myasthenia gravis.
    \item \textbf{Public health notification:} all suspected cases are reported so the food source can be investigated and removed.
\end{itemize}

\section*{Management}
Hospital care is essential, and the key treatment is antitoxin given as early as possible:
\begin{itemize}
    \item \textbf{Antitoxin:} neutralises circulating toxin and prevents further paralysis, but does not reverse weakness that has already developed; early administration is therefore critical.
    \item \textbf{Intensive care support:} mechanical ventilation for respiratory failure, tube feeding, and close monitoring of vital capacity.
    \item \textbf{Wound care:} for wound botulism, surgical debridement to remove infected tissue, which also reduces toxin production.
    \item \textbf{Antibiotics:} used in wound botulism; not recommended in uncomplicated infant botulism, where they can lyse bacteria and release more toxin.
    \item \textbf{Enemas and laxatives:} sometimes used in foodborne cases to eliminate toxin still in the gut, once swallowing is protected.
    \item \textbf{Rehabilitation:} physiotherapy and occupational therapy to rebuild strength and function during the long recovery.
    \item \textbf{Prevention of secondary complications:} physiotherapy, turning, and anticoagulation measures to prevent pneumonia, pressure sores, and blood clots while immobile.
\end{itemize}

\section*{Complications}
The effects of botulinum toxin can lead to serious complications:
\begin{itemize}
    \item Respiratory failure requiring a breathing machine, the most dangerous complication
    \item Aspiration pneumonia and other lung infections
    \item Prolonged generalised weakness and fatigue lasting weeks to months
    \item Pressure sores, deep vein thrombosis, and muscle wasting during immobility
    \item Urinary retention and constipation from autonomic involvement
    \item Long-term fatigue, breathlessness, and reduced exercise tolerance in some survivors
    \item Death in severe untreated cases; even with modern intensive care a small proportion die
\end{itemize}

\section*{Prognosis and Outcomes}
With prompt treatment, most people with botulism survive, but recovery is slow. Paralysis improves over weeks to months as new nerve endings grow to replace the damaged terminals, and full recovery commonly takes several months, sometimes up to a year. Around 5--10\% of foodborne cases are fatal even with modern intensive care, usually from breathing complications, and the figure is higher without antitoxin. Outcomes are substantially better when antitoxin is given early. Infants with botulism generally recover very well with supportive care alone, although the illness can be frightening for families and the hospital stay is often prolonged.

\section*{Prevention and Self-Care}
Botulism is almost entirely preventable with careful food handling and wound care:
\begin{itemize}
    \item Use safe, tested methods for home preserving, and pressure-cook low-acid foods such as vegetables and meats at the correct temperature.
    \item Boil home-preserved low-acid foods for the recommended time before eating, since boiling destroys the toxin.
    \item Never taste food from a bulging, leaking, or damaged can, or food with an off odour or appearance.
    \item Refrigerate garlic- or herb-infused oils and use them promptly.
    \item Do not give honey to babies under twelve months of age.
    \item Clean and dress wounds promptly, and seek medical care for deep, dirty, or puncture wounds.
    \item Discard suspect food rather than tasting it, and report suspected cases to health authorities.
\end{itemize}

\section*{Sources and Further Reading}
The following sources provide reliable, current information on botulism and its prevention.
\begin{itemize}
    \item NHS. \textit{Botulism.} https://www.nhs.uk/conditions/botulism/
    \item Mayo Clinic. \textit{Botulism: symptoms, causes, and treatment.} https://www.mayoclinic.org/
    \item Centers for Disease Control and Prevention. \textit{Botulism.} https://www.cdc.gov/botulism/
    \item World Health Organization. \textit{Botulism fact sheet.} https://www.who.int/
    \item MSD Manual. \textit{Botulism.} https://www.msdmanuals.com/
    \item MedlinePlus. \textit{Botulism.} https://medlineplus.gov/
\end{itemize}
